\documentclass[a4paper,12pt]{article}
\usepackage[utf8]{inputenc}
\usepackage[spanish]{babel}
\usepackage{amsmath}
\usepackage{geometry}

\geometry{margin=2.5cm}

\title{Solucionario: Examen de Física I - Forma C}
\date{}

\begin{document}
\maketitle

\section{Parte I: Teoría}
\begin{enumerate}
    \item \textbf{Definición}: Estudia materia, energía, tiempo y espacio.
    \item \textbf{Redondeo}: Menor cantidad de cifras significativas presentes.
    \item \textbf{Magnitudes SI}: Metro, Kilogramo, Segundo, Kelvin, Amperio.
    \item \textbf{Notación Científica}: $4.5 \times 10^{-7}$.
\end{enumerate}

\section{Parte II: Práctica}
\begin{enumerate}
    \item \textbf{Longitud}: $2.5 \times 100,000 / 2.54 = 98,425.2$ in.
    \item \textbf{Caudal}: $1.5 \times 3785.41 \times 3600 = 20,441,214$ cm³/hr.
    \item \textbf{Notación Científica}: $1.5 \times 10^{-8}$.
    \item \textbf{Cifras Significativas}: $12.45 \times 2.0 = 24.9 \rightarrow 25$.
\end{enumerate}

\end{document}
